%-*-latex-*-
%-*-latex-*-
\newcommand\COURSE{ciss350}
\newcommand\ASSESSMENT{q05}
\newcommand\ASSESSMENTTYPE{Quiz}
\newcommand\POINTS{\textwhite{xxx/xxx}}

\input{myquizpreamble}
\input{yliow}
\input{\COURSE}
%-*-latex-*-

%-*-latex-*-

%-*-latex-*-
\renewcommand\TITLE{\ASSESSMENTTYPE \ \ASSESSMENT}

\renewcommand\EMAIL{}
\renewcommand\AUTHOR{eagraham1@cougars.ccis.edu}

\textwidth=6in
\begin{document}
\topmatterthree

%------------------------------------------------------------------------------
\nextq
Complete the following recursive function to compute GCD.
(Yes I know this is straight from my notes. Do it anyway.)
\\
\ANSWER
\begin{answercode}
// Do NOT use for-loops or while-loops or do-while-loops
int gcd(int m, int n)
{
  if (n == 0)
  {
    return m;
  }
  else
  {
    return gcd(n, m % n);
  }
}
\end{answercode}

TURN PAGE

%------------------------------------------------------------------------------
\newpage
\nextq
Complete the following recursive function
\verb!print(m, n)! that prints a sequence
of consecutive integers starting with \verb!m!
and ending with \verb!n - 1!.
For instance the output of
\begin{console}[fontsize=\small]
print(2, 6);   // prints 2 3 4 5 and newline
print(10, 18); // prints 10 11 12 13 14 15 16 17 and newline
\end{console}
is
\begin{console}[fontsize=\small]
2 3 4 5
10 11 12 13 14 15 16 17
\end{console}
Note that the function must be recursive.
\\
\ANSWER
\begin{answercode}
// Do NOT use for-loops or while-loops or do-while-loops
void print(int m, int n)
{
    if (m >= n)
    {
        std::cout << ' ';        // <-- Write ONE statement here
    }
    else
    {
        std::cout << m << ' ';   // <-- Write ONE statement here
        print(m + 1, n);         // <-- Do not change this statement
    }
}
\end{answercode}

TURN PAGE

%------------------------------------------------------------------------------
\newpage
\nextq
Complete the following recursive function
\texttt{int count(int *start, int *end, int target)}
that counts the number of times \verb!target!
appears in
\verb!*start!, \verb!*(start + 1)!, \verb!*(start + 2)!,
..., \verb!*(end - 1)!.
For instance the output of
\begin{console}[fontsize=\small]
int x[] = {1, 2, 3, 4, 3, 2, 1, 2, 3, 4};
std::cout << count(x + 2, x + 8, 3) << '\n';
\end{console}
is \verb!2!.
Note that the function must be recursive.
\\
\ANSWER
\begin{answercode}
// Do NOT use for-loops or while-loops or do-while-loops
int count(int * start, int * end, int target)
{
    int c = 0;
    if (start >= end)
    {
        return 0;
    }
    else
    {
      if (*start == target)
      {
        c++;
      }
      return c + count(start + 1, end, target);
    }
}
\end{answercode}

%------------------------------------------------------------------------------
%-*-latex-*-
\newpage
\input{instructions}
\end{document}

